% TOP_PROBLEM: {"id": 30006971, "title": "The McKay-Navarro Conjecture for Finite Groups", "category_id": 4, "category": {"id": 4, "name": "algebra", "display_name": "Algebra", "description": "Group theory, ring theory, field theory, and algebraic structures.", "slug": "algebra", "order_index": 4, "created_at": "2026-07-31T15:26:25.671Z"}, "status": "open"}
% FIRST_POSTED: 2026-09-25
% ENTRY_KIND: statement_only
% !TeX program = lualatex
% Definition and sources only; this file is not an LLM solution attempt.
\documentclass[11pt]{article}
\usepackage[margin=25mm]{geometry}
\usepackage{fontspec}
\setmainfont{Latin Modern Roman}
\usepackage{amsmath,amssymb,mathtools}
\usepackage{unicode-math}
\setmathfont{Latin Modern Math}
\usepackage{xurl,hyperref}
\hypersetup{colorlinks=true,urlcolor=blue}
\setcounter{secnumdepth}{0}
\setlength{\parindent}{0pt}
\setlength{\parskip}{5pt}
\setlength{\emergencystretch}{3em}
\begin{document}
\section*{\texorpdfstring{The McKay-Navarro Conjecture for Finite Groups}{The McKay-Navarro Conjecture for Finite Groups}}\label{30006971}
\subsection{Definitions and mathematical statement}
Let $G$ be finite, $p$ prime, and $P$ a Sylow $p$-subgroup. Write $\operatorname{Irr}_{p'}(G)$ for its irreducible complex characters of degree not divisible by $p$. In $\mathcal G=\operatorname{Gal}({\mathbb{Q}}(\zeta_{|G|})/{\mathbb{Q}})$, let $\mathcal H_p$ consist of automorphisms whose action on roots of unity of order prime to $p$ is raising to a common $p$-power. It acts by $\chi^\sigma(g)=\sigma(\chi(g))$.

The standard conjecture stated in the cited paper asks for an $\mathcal H_p$-equivariant bijection
\[
 \operatorname{Irr}_{p'}(G)\longleftrightarrow
 \operatorname{Irr}_{p'}(N_G(P)).
\]
Equivariance means that applying any $\sigma\in\mathcal H_p$ before or after the bijection has the same effect. It does not require equivariance under arbitrary Galois automorphisms. The source proves the prime-two case and states the all-prime conjecture. The catalogue's extra word blockwise refers to a distinct refinement; no precise blockwise target is supplied there, and that refinement is not silently identified with this global bijection.
\par\smallskip\textit{Formulation and concept references:} \href{https://arxiv.org/abs/2211.14237}{[S1]}.

\subsection{Short English statement}
Can prime-to-p-degree characters of a finite group and its Sylow normalizer be matched compatibly with the specified Galois action?
\subsection{Sources}
\begin{itemize}
\item[S1]L. Ruhstorfer, A. A. Schaeffer Fry, arXiv:2211.14237 (2022); Adv. Math. 483 (2025), 110369.
\url{https://arxiv.org/abs/2211.14237}.
\textit{Formulation source. Consulted 24 September 2026: The full text states the H\_p-equivariant global bijection and distinguishes its blockwise refinement.}
\end{itemize}
\end{document}
